\section{Adaptation on full affine Hilbert balls}
\label{app:hilbert-adaptation}

The geometry in this appendix allows overlapping features and infinite
rank.  Let $B_{\mathcal H}$ denote the closed unit ball of a real RKHS.
Assume that, for $\nu\in\{\mu,\pi\}$,
\begin{equation}
 \cG_\nu=c_\nu+R_\nu B_{\mathcal H_\nu},\qquad R_\nu\ge0,
 \label{eq:hilbert-full-balls}
\end{equation}
with the pointwise envelope and supplied sequential-M1 cores of
Appendix~\ref{app:model}.  These are full affine balls; the two RKHSs and
their centers can differ.  A zero radius gives a singleton.  The centers,
radii, and kernels specify the geometric subclass, but need not be
additional inputs to the estimator.

\begin{theorem}[Adaptation on full affine Hilbert balls]
\label{thm:hilbert-adaptation}
For every represented pair satisfying \eqref{eq:hilbert-full-balls},
there is one Borel rule $\delta_{\Gamma,\rho,n}:\mathcal Z^{2n}\to[-3,3]$
such that, simultaneously for all $(a,s,b,t)\in[0,1]^4$ and every
$P\in\mathcal P_{\rho,n,a,s,b,t}(\Gamma)$,
\begin{equation}
 \E_P|\delta_{\Gamma,\rho,n}-\beta(P)|
 \le C_\rho(q+ab)\le C_\rho(q+W).
 \label{eq:hilbert-adaptation-rate}
\end{equation}
The rule uses none of the four certificates.  Its bound is independent of
the dimensions, ranks, and overlap of the two feature spaces.  In fact,
the upper bound does not use the localized-complexity inequalities.
\end{theorem}

\begin{proof}
Write $N=2n$, $f_X=(f(X_1),\ldots,f(X_N))'$, and
$\|z\|_N^2=z'z/N$.  First recover the geometry from the supplied cores.
For each side set
\begin{equation}
 \begin{aligned}
 \widetilde c_\nu(x)
  &=\tfrac12\{\sup_j g_{\nu,j}(x)+\inf_j g_{\nu,j}(x)\},\qquad
  z_{\nu,j}=g_{\nu,j}-\widetilde c_\nu,\\
 u_{\nu,\pm}(x,y)&=\sup_j\{z_{\nu,j}(x)\pm z_{\nu,j}(y)\},\qquad
 \ell_\nu(x,y)=\tfrac14\{u_{\nu,+}(x,y)^2-u_{\nu,-}(x,y)^2\}.
 \end{aligned}
 \label{eq:hilbert-core-recovery}
\end{equation}
Every extremum is finite and Borel.  Pointwise sequential M1 preserves
the supremum of every finite linear combination of evaluations.  For a
kernel $k_\nu$ representing $\mathcal H_\nu$, symmetry of its ball gives
\[
 \widetilde c_\nu=c_\nu,\qquad
 u_{\nu,\pm}(x,y)
   =R_\nu\|k_\nu(x,\cdot)\pm k_\nu(y,\cdot)\|_{\mathcal H_\nu}.
\]
Polarization therefore shows that
\begin{equation}
 \ell_\nu(x,y)=R_\nu^2k_\nu(x,y),\qquad
 |c_\nu(x)|+\sqrt{\ell_\nu(x,x)}\le B_{\mathcal G}.
 \label{eq:hilbert-kernel-envelope}
\end{equation}
These statements also hold when $R_\nu=0$.  Thus the recovered $\ell_\nu$
is a Borel positive semidefinite kernel, whether or not a kernel was
supplied separately.

If $N<16B_0^2/v_-$, define $\delta=0$.  Its loss is at most three,
which is absorbed by $C_\rho q$.  Otherwise form the matrices
\begin{equation}
 \begin{aligned}
 H_\nu&=[\ell_\nu(X_i,X_j)/N]_{i,j},\qquad
 C_X=[c_{\mu,X},c_{\pi,X}],\\
 M_X&=I_N-C_X(C_X'C_X)^\dagger C_X',\qquad
 H=M_X(H_\mu+H_\pi)M_X,\\
 \lambda_N&=\frac{16B_{\mathcal G}^2B_0^2}{Nv_-},\qquad
 A_X=M_X(I_N+H/\lambda_N)^{-1}M_X.
 \end{aligned}
 \label{eq:hilbert-contraction}
\end{equation}
Here $\dagger$ denotes the Moore--Penrose inverse.  The center projection
is Borel, for example by taking the limit of
$C_X(C_X'C_X+k^{-1}I_2)^{-1}C_X'$ as $k\to\infty$.
All other operations are Borel and the displayed ordinary inverse always
exists.  For $D_A=T'A_XT/N$, the rule is
\begin{equation}
 \delta=\Pi_{[-3,3]}
 \left\{\frac{T'A_XY/N}{D_A\vee(v_-/4)}\right\}.
 \label{eq:hilbert-rule}
\end{equation}
It is defined on every sample, including singular designs and zero
kernels.  In particular, it does not estimate the approximation errors.

Since $H=M_XHM_X$, the matrices $M_X$ and $H$ commute.  Hence
$A_X$ is symmetric, $0\preceq A_X\preceq I_N$, and it annihilates both
center vectors.  The envelope in \eqref{eq:hilbert-kernel-envelope} gives
$\operatorname{tr}H\le2B_{\mathcal G}^2$, so, on every sample,
\begin{equation}
 \begin{aligned}
 \operatorname{tr}(I_N-A_X)
 &=\operatorname{rank}(I_N-M_X)
   +\operatorname{tr}\{H(H+\lambda_N I_N)^{-1}\}\\
 &\le2+2B_{\mathcal G}^2/\lambda_N
 \le Nv_-/(4B_0^2).
 \end{aligned}
 \label{eq:hilbert-trace}
\end{equation}
For either side,
\begin{equation}
 A_XH_\nu A_X
 =A_XM_XH_\nu M_XA_X
 \preceq A_XHA_X\preceq(\lambda_N/4)I_N,
 \label{eq:hilbert-residual-matrix}
\end{equation}
where the last inequality is $x/(1+x)^2\le1/4$ for $x\ge0$.
The evaluation operator $f\mapsto R_\nu f_X/\sqrt N$ has Gram matrix
$H_\nu$.  Because $A_Xc_{\nu,X}=0$, it follows that, with
$e_N=\sqrt{\lambda_N}/2$,
\begin{equation}
 \sup_{g\in\cG_\mu}\|A_Xg_X\|_N\le e_N,\qquad
 \sup_{h\in\cG_\pi}\|A_Xh_X\|_N\le e_N.
 \label{eq:hilbert-comparator-residual}
\end{equation}
No population covariance inverse or eigenvalue-decay condition enters
these inequalities.

For the risk calculation, first take fixed full-class comparators $g,h$
with $\|\mu_0-g\|_2\le a$ and $\|\pi_0-h\|_2\le b$.
Write $r_\mu=\mu_0-g$, $r_\pi=\pi_0-h$ and suppress the subscript on
$A_X$.  Expanding $\pi_{0,X}'A\mu_{0,X}$ into the four comparator and
remainder terms, empirical Cauchy--Schwarz followed by expectation
Cauchy--Schwarz and \eqref{eq:hilbert-comparator-residual} gives
\begin{equation}
 \E|\pi_{0,X}'A\mu_{0,X}/N|
 \le ab+(a+b)e_N+B_{\mathcal G}e_N.
 \label{eq:hilbert-bias}
\end{equation}
For example, the public--public term is at most
$\|h_X\|_N\|Ag_X\|_N\le B_{\mathcal G}e_N$; the
remainder--remainder term has expectation at most $ab$.

Conditional on $X_{1:N}$, the $u_i$ are independent and centered.
Conditional on $(X_{1:N},T_{1:N})$, the $\varepsilon_i$ are likewise
independent and centered.  Boundedness and contraction therefore imply
\[
 \E\{(u'A\mu_{0,X}/N)^2\mid X_{1:N}\}\le B_0^4/N,
 \qquad
 \E\{(T'A\varepsilon/N)^2\mid X_{1:N},T_{1:N}\}\le B_0^4/N.
\]
Thus $Q_A=T'A(Y-\beta_0T)/N$ satisfies
\begin{equation}
 \E|Q_A|
 \le ab+(a+b+B_{\mathcal G})e_N+2B_0^2/\sqrt N.
 \label{eq:hilbert-numerator}
\end{equation}

It remains to control the denominator under only average residual
variance.  Put $v_i=\E(u_i^2\mid X_i)$,
$\overline v=N^{-1}\sum_i v_i$, and $\Omega=\operatorname{diag}(v_i)$.
Since $0\le v_i\le B_0^2$, \eqref{eq:hilbert-trace} yields
\[
 \E(D_A\mid X_{1:N})
 =\{\pi_{0,X}'A\pi_{0,X}+\operatorname{tr}(A\Omega)\}/N
 \ge\overline v-v_-/4.
\]
Moreover $\E\overline v=v\ge v_-$, so
$\Pp(\overline v<3v_-/4)\le16B_0^4/(Nv_-^2)$.
Conditional independence and centering give
\[
 \begin{aligned}
 \operatorname{Var}(u'Au\mid X_{1:N})
 &=\sum_i A_{ii}^2\operatorname{Var}(u_i^2\mid X_i)
   +4\sum_{i<j}A_{ij}^2v_iv_j
 \le2B_0^4\operatorname{tr}(A^2)\le2B_0^4N,\\
 \operatorname{Var}(2u'A\pi_{0,X}\mid X_{1:N})
 &\le4B_0^4N.
 \end{aligned}
\]
Using $\operatorname{Var}(U+V)\le2\operatorname{Var}U+2\operatorname{Var}V$
therefore gives $\operatorname{Var}(D_A\mid X_{1:N})\le12B_0^4/N$;
no assumption about third moments being zero is needed.  On
$\{\overline v\ge3v_-/4\}$ the conditional mean is at least $v_-/2$.
Another Chebyshev bound consequently proves
\begin{equation}
 \Pp(D_A<v_-/4)\le208B_0^4/(Nv_-^2).
 \label{eq:hilbert-denominator}
\end{equation}
On the complementary event the floor is inactive and clipping cannot
increase the error, so $|\delta-\beta_0|\le4|Q_A|/v_-$.
On the bad event the clipped loss is at most six.  Combining
\eqref{eq:hilbert-numerator}--\eqref{eq:hilbert-denominator} gives the
explicit finite-sample bound
\begin{equation}
 \E|\delta-\beta_0|
 \le\frac4{v_-}
 \left\{ab+(a+b+B_{\mathcal G})e_N+\frac{2B_0^2}{\sqrt N}\right\}
 +\frac{1248B_0^4}{Nv_-^2}.
 \label{eq:hilbert-explicit-risk}
\end{equation}
For unattained approximation infima, use fixed comparators with errors
$a+\eta,b+\eta$ throughout and let $\eta\downarrow0$; neither comparator
is an algorithm input.  Finally $e_N=2B_{\mathcal G}B_0/\sqrt{Nv_-}$,
$a,b\le1$, and $ab\le W$.  This proves
\eqref{eq:hilbert-adaptation-rate}, including every zero-budget face.
\end{proof}

For a finite implementation, let the centers $c_\nu$ and features
$\phi_\nu:\mathcal X\to\mathbb R^{d_\nu}$ be explicitly evaluable Borel
functions, and let $Q_\nu$ be a public positive semidefinite matrix.  The full ellipsoid
\[
 \cG_\nu=c_\nu+
 \{\phi_\nu'Q_\nu^{1/2}\theta:\|\theta\|_2\le1\}
\]
has kernel $\ell_\nu(x,y)=\phi_\nu(x)'Q_\nu\phi_\nu(y)$, and its
envelope condition is exactly
$|c_\nu(x)|+\sqrt{\phi_\nu(x)'Q_\nu\phi_\nu(x)}\le B_{\mathcal G}$.
The rule then uses finite matrix operations, even if $d_\nu>N$ or the
features are linearly dependent.  Arbitrary overlap within and between
the two feature maps is permitted.  Bounded-diagonal RKHS balls give the
corresponding infinite-rank examples, without any spectral-decay
assumption.  The same upper proof also applies to classes merely
\emph{contained} in such balls if the enclosing centers and kernels,
with the displayed diagonal bound, are supplied as additional public
geometry; the recovery in \eqref{eq:hilbert-core-recovery} is asserted
only for the full balls.  An arbitrary bounded $\ell_1$ dictionary does
not automatically have this Hilbert-ball envelope.

The construction uses classical penalized partialling-out: the centers
are unpenalized nuisance directions and $H$ supplies the remaining ridge
smoother.  Partial smoothing splines and partially linear kernel ridge
regression are established methods
\citep{heckman1986spline,zhao2016partially}; series residualization also
uses product approximation-bias arguments
\citep{cattaneo2018alternative}.  The result here is a bounded
finite-sample specialization with unknown approximation errors and
recovery from the represented full balls; no novelty of the underlying
regression method is claimed.

For completeness, restrict the outer public-pair supremum in
\eqref{eq:risk} to the full affine Hilbert balls above, including radius
zero.  On the strict lower-endpoint slack stratum $0<v_-<B_0^2$, its
minimax risk is $\asymp_\rho q+ab$.  The upper bound is
Theorem~\ref{thm:hilbert-adaptation}.  For the lower bound,
Appendices~\ref{app:q-lower} and~\ref{app:ab} both use the identical
singleton pair $\cG_\mu=\cG_\pi=\{0\}$, with zero difference
complexities and repeated-zero M1 cores.  Their bounds (C.9) and (E.13)
therefore hold in this restricted outer class for every $s,t$;
$\max(c_{q,\rho}q,c_{ab,\rho}ab)\ge
\tfrac12\min(c_{q,\rho},c_{ab,\rho})(q+ab)$ gives the claim.
This is an outer-class lower bound, not a lower bound for each prescribed
nondegenerate Hilbert-ball pair.

Budget-free adaptation is also an established objective: SADE
Theorem 2.1 and Remark 2.1 of \citet{gu2026tamingbiases} give a
fixed-outcome-class high-probability bound of order
\(\sqrt{\tau/n}+a+s^2+\lambda\), without those budget inputs.
At fixed tail parameter and \(\lambda\lesssim q\), its leading scale
is \(q+a+s^2\). This one-class guarantee is different from the
two-class expected-\(L_1\) result \(q+ab\) here. Neither comparison claims a new general
adaptation principle or an exhaustive exclusion of earlier special cases.
